\section{Cavity increments and attainment of the entropy dual}
\label{cav:section}

Adding finitely many coordinates follows the cavity variational method of
Aizenman, Sims, and Starr~\cite{aizenmansimsstarr2003}. For the perceptron,
the new coordinates first produce the gradient covariance identified in
Section~\ref{bulk:section}. The purpose of this section is to evaluate that
increment and prove that its field attains the Ising entropy dual.

Fix an integer $L\ge1$. All limits and error terms in this section keep
$L$ fixed until the final limit after telescoping. For every deterministic
$v\in[1,2]^{\mathbb N}$ define the unnormalized increment
\begin{equation}\label{cav:increment-definition}
 \Delta_{N,L}(v)=\mathbb E\log Z_{N+L}^v-\mathbb E\log Z_N^v,
 \qquad d=d_{N,L}=M_{N+L}-M_N.
\end{equation}
The same infinite parameter sequence is used in both dimensions.
Notice that $d$ is bounded at fixed $L$ and $|d-\alpha L|\le1$.

\subsection{A uniform finite-size cavity comparison}

Write a spin in dimension $N+L$ as $(x,z)$ with prior
$\nu_N\otimes\nu_L$, and split each pattern vector as $(g^a,y_a)$.
The blocks are independent standard Gaussians, with $y_a\in\mathbb R^L$.
The larger pattern argument is
\[
 \lambda_N b_a(x)+\frac{y_a\cdot z}{\sqrt{N+L}},
 \qquad \lambda_N=\sqrt{\frac{N}{N+L}},
\]
where $b_a(x)=g^a\cdot x/\sqrt N$, also for the $d$ new patterns.

First replace the larger system's perturbation by exactly the bulk
perturbation $H_N^v(x)$.  To justify this replacement, write
\[
 R^+((x,z),(x',z'))=\frac{NR(x,x')+z\cdot z'}{N+L}.
\]
Then $|R^+-R|\le2L/N$.  The series
$\xi_N^v$ and their first derivatives are bounded uniformly on
$[-1,1]$, and the omitted series tail is $O(4^{-N})$.  Consequently
\begin{align}
 &\sup_{x,x',z,z',v}
 \left|s_{N+L}^2\xi_{N+L}^v(R^+)-s_N^2\xi_N^v(R)\right|
 \nonumber\\
 &\qquad\le C_L\left(N^{-1/4}+N^{3/4}4^{-N}\right)=o(1).
 \label{cav:perturbation-comparison}
\end{align}
Lemma~\ref{arr:calculus}, conditionally on all patterns, controls by
this bound both the expected unnormalized log partition function and
the disorder-averaged Gibbs expectation of any fixed bounded
finite-replica spin test independent of the perturbing fields.

Conditionally on the full size-$N$ bulk data, introduce the vector
Gaussian field
\begin{equation}\label{cav:linear-field}
 \mathcal Y(x)=\frac1{\sqrt N}\sum_{a=1}^{M_N} f'(b_a(x))y_a.
\end{equation}
Its $L$ coordinates are conditionally independent centered Gaussian
fields of covariance $B_N$, independent also of the $d$ new bulk
pattern fields.  Define
\begin{align}
 W(x,z)&=\exp\left\{z\cdot\mathcal Y(x)+\frac L2C_N(x)
                   +\sum_{a=M_N+1}^{M_N+d}f(b_a(x))\right\},
 \label{cav:W}\\
 Q(x,z)&=\frac1{2N}\sum_{a=1}^{M_N} f''(b_a(x))
                      \bigl((y_a\cdot z)^2-L\bigr).
 \label{cav:Q}
\end{align}
Here $B_N,C_N$ are the full bulk observables in
\eqref{bulk:B}--\eqref{bulk:C}.

\begin{lemma}[Cavity comparison for increments and replica laws]
\label{cav:finite-comparison}
Let
\begin{equation}\label{cav:untilted-base}
 G_0=G_N^v\otimes\nu_L
\end{equation}
be the size-$N$ bulk Gibbs probability times the independent cavity-spin
prior, \emph{before} multiplication by $W$.  There is a deterministic
$\kappa_L>0$ such that $G_0(W)\ge\kappa_L$ for every realization and
all sufficiently large $N$.  Uniformly over all allowed $v$,
\begin{equation}\label{cav:finite-increment}
 \Delta_{N,L}(v)=\mathbb E\log G_0(W)+o(1).
\end{equation}
The size-$(N+L)$ Gibbs expectations of any fixed bounded
finite-replica spin test are likewise approximated, after the above
coupling, by the measure $W\,dG_0/G_0(W)$, with disorder-averaged error
$o(1)$ times the test's bound.  In particular this applies to bounded
continuous functions of $R^+_{12}$ and their products with $z_1^1z_1^2$;
inside such continuous tests one may also replace $R^+_{12}$ by $R_{12}$.
Finally, $\Delta_{N,L}(v)\ge-C_L$ uniformly for large $N$.
\end{lemma}

\begin{proof}
Compact support makes all scale derivatives needed below bounded.
Uniformly in $s\in\mathbb R$, Taylor expansion in the scale parameter
therefore gives
\begin{align*}
 f(\lambda_Ns)&=f(s)-\frac L{2N}s f'(s)+O_L(N^{-2}),\\
 f^{(j)}(\lambda_Ns)&=f^{(j)}(s)+O_L(N^{-1}),\qquad j=1,2.
\end{align*}
For example, derivatives of $f(\lambda s)$ with respect to $\lambda$
are $s^r f^{(r)}(\lambda s)$, uniformly bounded when $\lambda$ stays
near one.  Next expand the translation by
$(y_a\cdot z)/\sqrt{N+L}$ through second order, and replace its
denominators by $\sqrt N$ and $N$.  Its third-order remainder is bounded
by $C_LN^{-3/2}|y_a|^3$, uniformly over $x,z$.  The scale and denominator
replacements add terms bounded by
$C_L(N^{-2}+N^{-3/2}|y_a|+N^{-2}|y_a|^2)$.
Summing over the $M_N=O(N)$ old rows thus gives an expected uniform
error $O_L(N^{-1/2})$.

The deterministic part of the quadratic term is
$L\sum_a f''(b_a(x))/(2N)$; together with the dilation term it is
$LC_N(x)/2$.  The remaining quadratic term is precisely $Q$, and the
linear term is $z\cdot\mathcal Y(x)$.  Each new row can be replaced by
$f(b_a(x))$: its expected uniform cost is at most
$C_L(N^{-1}+N^{-1/2}\mathbb E|y_a|)$.  Since there are only $d=O_L(1)$
new rows, this is $o(1)$ in total.  We have proved that, after the
perturbation comparison, the energy equals
\begin{equation}\label{cav:taylor}
 \sum_{a=1}^{M_N}f(b_a(x))+H_N^v(x)
       +\log W(x,z)+Q(x,z)+\operatorname{rem}_N(x,z),
 \qquad \mathbb E\sup_{x,z}|\operatorname{rem}_N(x,z)|=o(1).
\end{equation}
The elementary log-partition bound by the uniform energy difference
removes this remainder.  Differentiating a normalized $r$-replica
expectation along the energy segment bounds its change by
$2r\|D\|_\infty\sup|\operatorname{rem}_N|$, so the same removal applies
to the asserted tests.

Put $A(x)=LC_N(x)/2+\sum_{a>M_N}f(b_a(x))$.  Uniformly,
$|A(x)|\le C_L$.  Since the cavity prior is a product of symmetric
signs,
\[
 G_0(W)=G_N^v\left(e^{A(x)}
                    \prod_{i=1}^L\cosh(\mathcal Y_i(x))\right)
 \ge e^{-C_L}=:\kappa_L.
\]
This bound concerns the \emph{untilted} probability
\eqref{cav:untilted-base}; it will control all subsequent
normalizations.

It remains to remove $Q$, which need not be uniformly small over all
spin states.  For each fixed $x,z$ and fixed bulk data, the variables
$y_a\cdot z$ are independent $N(0,L)$ random variables.  Thus
\begin{equation}\label{cav:Q-moments}
 \mathbb E_y Q(x,z)^2\le\frac{C_L}{N},\qquad
 \mathbb E_y e^{t|Q(x,z)|}\le C_{L,t}\quad(t>0),
\end{equation}
uniformly for large $N$.  To check the second assertion directly, put
$c_a=f''(b_a(x))/(2N)$.  For fixed $t$ and sufficiently large $N$,
\[
 \log\mathbb E_y e^{t c_a((y_a\cdot z)^2-L)}
 =-tLc_a-\tfrac12\log(1-2tLc_a)=O_{L,t}(c_a^2).
\]
The linear terms cancel.  Summing over $a$ is $O_{L,t}(N^{-1})$;
using $e^{t|Q|}\le e^{tQ}+e^{-tQ}$ proves the claim.  Also
\begin{equation}\label{cav:linear-moments}
 \mathbb E_y e^{t z\cdot\mathcal Y(x)}\le C_{L,t},
\end{equation}
because its conditional variance is $LB_N(x,x)\le C_L$.

Write $P_W=W\,dG_0/G_0(W)$ for the normalized approximation.  The
bounds just proved imply
\begin{equation}\label{cav:Q-under-tilt}
 \mathbb E P_W(|Q|)+\mathbb E P_W(|e^Q-1|)
 \le C_LN^{-1/2}.
\end{equation}
Indeed $G_0(W)\ge\kappa_L$ and the non-linear-field part of $W$ is
bounded.  For the second term, use $|e^Q-1|\le|Q|e^{|Q|}$ and
Cauchy--Schwarz to bound its fixed-state $y$ expectation by
\[
 (\mathbb E_yQ^2)^{1/2}
 \bigl(\mathbb E_y e^{2z\cdot\mathcal Y(x)+2|Q|}\bigr)^{1/2}
 \le C_LN^{-1/2}.
\]
A further Cauchy--Schwarz inequality controls the last exponential by
\eqref{cav:Q-moments} and \eqref{cav:linear-moments}.  This uses no
independence between $Q$ and $\mathcal Y$.

Jensen's inequality and $\log t\le t-1$ now give
\[
 -P_W(|Q|)\le\log P_W(e^Q)\le P_W(|e^Q-1|),
\]
so removing $Q$ changes the expected log partition by $o(1)$.
For normalized tests, the following elementary bound avoids any
additional lower bound for $P_W(e^Q)$.  For every probability $P$ and
positive integrable $Y$, writing $m=P(Y)>0$,
\begin{equation}\label{cav:normalized-comparison}
 P\left|\frac Ym-1\right|
 \le P\left|\frac Ym-Y\right|+P|Y-1|
 =|1-m|+P|Y-1|\le2P|Y-1|.
\end{equation}
Apply this with $P=P_W,Y=e^Q$ and then use the product-measure bound
for a fixed number of replicas.  Together with
\eqref{cav:perturbation-comparison} and \eqref{cav:taylor}, this proves
\eqref{cav:finite-increment} and all the claimed sample comparisons.
Uniform continuity on $[-1,1]$ handles the replacement of $R^+$ by $R$.
The deterministic denominator lower bound proves the final uniform
lower bound for the increments.
\end{proof}

\subsection{Evaluation of an arbitrary good subsequence}

Take an arbitrary deterministic good parameter sequence, possibly
already along a subsequence of dimensions.  Extract further as in
Proposition~\ref{bulk:restoration}, and also so that $d_{N,L}$ is
constant, say equal to $d$.  Let $p$ be the nondecreasing quantile of
$\mu$ and put
\begin{equation}\label{cav:field-path}
 h(u)=a(p(u)),\qquad 0<u<1.
\end{equation}
This is a bounded nonnegative nondecreasing path, and its values are
measurable with respect to the completed sigma-algebra generated by
$p$.  Values on null sets may be chosen arbitrarily.  The limiting
conditional field covariances in $W$ are $a(R_{lm})$ off the diagonal
and $a_{\mathrm d}$ on the diagonal for each linear-field component,
and $R_{lm}$ with diagonal one for each new pattern.  The remaining
base observable $C_N$ converges to the constant $c$.

The covariance and cavity-correction limits will give the increment
$L\{\ell(h)+\frac12\int_0^1ph\}+dV_f(p)$.
The remaining issue is whether its field contribution attains
$S_{\mathrm I}(p)$, rather than merely being bounded above by it.
We first evaluate the increment by finite ordered approximations of
$(p,h)$.  Coordinate exchangeability will then show that the derivative
paths of these finite field recursions converge to $p$.  This is the
identity that turns the supporting inequality of
Theorem~\ref{var:concavity} into entropy attainment.

To justify passage through the overlap limit despite the unbounded
linear field, let $[t]_D=\max(-D,\min(t,D))$ and define
\[
 W_D(x,z)=\exp\{[z\cdot\mathcal Y(x)]_D+A(x)\}.
\]
Symmetry under $z\mapsto-z$, and oddness of clipping, give
$G_0(W_D)\ge\kappa_L$ with the same constant as before.
For a centered Gaussian $T$ of variance at most $C_L$,
\[
 \sup_{\operatorname{Var}(T)\le C_L}
 \mathbb E\big[e^{|T|}\mathbf 1_{\{|T|>D\}}\big]\longrightarrow0.
\]
For example, on $|T|>D$ the integrand is at most
$e^{-D}e^{2|T|}$, whose expectation is uniformly bounded.  Hence
\begin{equation}\label{cav:truncation-error}
 \sup_{N,v}\mathbb E G_0(|W-W_D|)\longrightarrow0
 \qquad(D\to\infty).
\end{equation}
The denominator bounds control the expected log difference by
\[
 \kappa_L^{-1}\mathbb E G_0(|W-W_D|).
\]
The $L^1$ distance between the normalized measures is at most
$2\kappa_L^{-1}G_0(|W-W_D|)$, and the analogous bound with a factor
$r$ applies to $r$ replicas.  Thus bounded replica tests have the same
uniform truncation control.

For fixed $D$, the factor $W_D$ itself is bounded above and away from
zero and is continuous in the finite field evaluations and in $C_N$.
Lemma~\ref{arr:limit} applies: finite-dimensional Gaussian covariance
continuity, followed by polynomial approximation of reciprocal
normalizers and logarithms, expresses the limit through finitely many
base replicas.  This applies jointly to all field components and the
bounded continuous overlap tests under consideration.  The same
truncation estimate holds for each limiting finite cascade because its
true diagonal field variance is still $a_{\mathrm d}$.  Clipping is
performed on the full field, including its residual, before any tilt.
We may therefore evaluate unclipped finite cascade expressions and
then pass to their limit, by first fixing $D$, taking the array and
rounding limits, and finally sending $D$ to infinity.

To make the latter calculation precise, bin the ordered pair
$(r,a(r))$ jointly on successively finer meshes of
$[0,1]\times[0,a_{\mathrm d}]$.  Retain the bins of positive
$\mu$-probability and choose representatives ordered in both
coordinates.  The resulting step paths will be denoted
$(\bar p_m,\bar h_m)$.  They converge to $(p,h)$ in $L^1$, and their
values approximate both coordinates on the sampled off-diagonal
entries.  Leave all diagonal variances equal to $1$ and
$a_{\mathrm d}$.  This construction is a finite nondecreasing
function of $r$, so Lemmas~\ref{arr:quantization} and~\ref{arr:cascade}
apply.  It never divides an interval on which $p$ is constant.
This last restriction is important for entropy attainment: tests of
the spin overlap identify only a conditional mean of the field
derivative.  Keeping the derivative paths measurable as functions of
$p$ will let those tests identify the entire limiting path.

Write $\bar h_{m,\mathrm{top}}$ for the last step value of $\bar h_m$.
In such a finite cascade, the linear-field residual on each visit has
coordinate variance $a_{\mathrm d}-\bar h_{m,\mathrm{top}}$.
Averaging its exponential against a spin $z$ contributes the constant
\[
 \frac L2(a_{\mathrm d}-\bar h_{m,\mathrm{top}})
\]
to the leaf log factor, since $|z|^2=L$.  Averaging over the cavity
spins then gives a sum of $L$ independent $\log\cosh$ factors for the
shared fields.  The field recursion of Definition~\ref{def:recursions}
subtracts $\bar h_{m,\mathrm{top}}/2$, so these shared factors contribute
$L(\ell(\bar h_m)+\bar h_{m,\mathrm{top}}/2)$.
Each of the $d$ independent pattern components contributes
$V_f(\bar p_m)$, including its final coefficient-one residual average.
The deterministic bulk factor contributes $Lc/2$.
The additive-component assertion of Lemma~\ref{arr:tilt} therefore
evaluates the expected log integral as
\[
 L\left(\ell(\bar h_m)+\frac{c+a_{\mathrm d}}2\right)
        +dV_f(\bar p_m).
\]
The continuity established in Section~\ref{var:section}, together with
\eqref{bulk:ibp-balance}, proves the following exact subsequential
limit:
\begin{equation}\label{cav:increment-limit}
 \lim_N\Delta_{N,L}(v(N))
 =L\left(\ell(h)+\frac12\int_0^1p(u)h(u)\,du\right)+dV_f(p).
\end{equation}

\subsection{Coordinate exchangeability and the conjugate field}

The identity established here is the perceptron analogue of the
critical relation identifying an overlap path with the derivative of
the one-spin cascade transform. Such relations are central to Chen
and Mourrat~\cite[Propositions~7.2 and~7.3]{chenmourrat2025} and
Chen, Issa, and Mourrat~\cite[Theorem~6.1]{chenissamourrat2026}. We do not
apply their Gaussian spin-glass theorems to this nonlinear Hamiltonian:
the restored covariance and coordinate exchangeability supply the
identity here. Once it is established, the centered-Ising concavity
credited in Theorem~\ref{var:concavity} gives the supporting inequality
that proves attainment of the dual supremum.

\begin{proposition}[The cavity field attains the entropy dual]
\label{cav:dual-attainment}
For every subsequential pair $(p,h)$ just constructed,
\begin{equation}\label{cav:entropy-identity}
 S_{\mathrm I}(p)=\ell(h)+\frac12\int_0^1p(u)h(u)\,du.
\end{equation}
The right side uses the continuous extension of $\ell$ to bounded
nondecreasing fields.
\end{proposition}

\begin{proof}
We identify the field-recursion derivative using one distinguished
cavity coordinate.  For a fixed rounding $(\bar p_m,\bar h_m)$, let
$\rho_{\bar h_m}$ be the bounded nondecreasing derivative path of
Section~\ref{var:section}; on level $i$ it equals
$\mathbb E[U_i'(Z_i)^2]$ for the field recursion.

Under the cavity tilt, the conditional mean of $z_1^1z_1^2$ given
pair level $i$ is precisely $(\rho_{\bar h_m})_i$.  To verify this
statement including the residual noise, fix the shared leaf field
$t$ for coordinate one.  If $\eta$ is the independent residual with
variance $a_{\mathrm d}-\bar h_{m,\mathrm{top}}$, then
\[
 \mathbb E_\eta e^{z_1(t+\eta)}
   =e^{z_1t}e^{(a_{\mathrm d}-\bar h_{m,\mathrm{top}})/2}.
\]
The second factor is independent of $z_1$, so after residual averaging
the spin magnetization is $\tanh t$.  Two visits, even to the same
leaf, use separate residuals and separate spin draws.  Below their
shared level their mark paths are conditionally independent.
Propagating the leaf magnetization backwards along the tilted kernels
therefore gives the conditional mean $U_i'$ at level $i$, by the
spatial derivative martingale for the field recursion.  The conditional
product mean is its averaged square.  All other coordinates and all
new pattern marks have separate kernels by the additive-component
property.  The shape law itself is unchanged by the tilt.

At finite size the disorder law of the original $(N+L)$-spin model,
including its perturbation at fixed $v$, is invariant under all
coordinate permutations.  Thus for every bounded continuous
$D_0:[-1,1]\to\mathbb R$,
\begin{equation}\label{cav:coordinate-identity}
 \mathbb E\langle D_0(R^+_{12})z_1^1z_1^2\rangle_{N+L}
 =\mathbb E\langle D_0(R^+_{12})R^+_{12}\rangle_{N+L}.
\end{equation}
Indeed the left side is unchanged if the distinguished coordinate is
replaced by any of the $N+L$ coordinates, and their average product is
$R^+_{12}$.  This symmetry is used before the cavity approximations.
Lemma~\ref{cav:finite-comparison}, the truncation passage above, and the
finite cascade calculation then imply
\begin{equation}\label{cav:rho-test}
 \lim_{m\to\infty}\int_0^1
    D_0(\bar p_m(u))\rho_{\bar h_m}(u)\,du
 =\int_0^1 D_0(p(u))p(u)\,du.
\end{equation}
For the right side, the unclipped cascade tilt preserves the rounded
overlap law by Lemma~\ref{arr:tilt}.  The truncation estimates above
transfer this evaluation to the finite-size limit, whose overlap law
has quantile $p$.

The paths $\rho_{\bar h_m}$ take values in $[0,1]$ and are
nondecreasing.  Bounded monotone selection and dominated convergence
give a subsequence converging in $L^1(0,1)$ to a function $\rho$.
Every $\rho_{\bar h_m}$ is measurable with respect to the same
completed sigma-algebra $\sigma(p)$: the bins were formed from
$(p,a(p))$ and hence are functions of $p$, with no splitting of its
flat intervals.  The subspace of $\sigma(p)$-measurable functions is
closed in $L^1(0,1)$.  Explicitly, conditional expectation on the Lebesgue probability space
$(0,1)$ onto $\sigma(p)$ is an $L^1$ contraction, so
\[
 \|\rho-\mathbb E[\rho\mid\sigma(p)]\|_1
 \le2\|\rho-\rho_{\bar h_m}\|_1\longrightarrow0.
\]
Thus $\rho$ itself is $\sigma(p)$-measurable.

Pass to this subsequence in \eqref{cav:rho-test}.  Uniform continuity of
$D_0$ and the $L^1$ convergence of $\bar p_m$ justify replacing
$D_0(\bar p_m)$ by $D_0(p)$.  We obtain
\[
 \int_0^1D_0(p(u))(\rho(u)-p(u))\,du=0
 \quad\hbox{for every continuous }D_0.
\]
Continuous tests determine finite signed measures on $[0,1]$.
Therefore the conditional expectation of $\rho-p$ given $p$ is zero.
Since $\rho-p$ is already $\sigma(p)$-measurable, it follows that
$\rho=p$ almost everywhere.  In particular, along this subsequence,
\begin{equation}\label{cav:rho-convergence}
 \rho_{\bar h_m}\longrightarrow p\quad\hbox{in }L^1(0,1).
\end{equation}

Apply the supporting inequality \eqref{var:support} with base
$\bar h_m$ and an arbitrary fixed nonnegative nondecreasing field step
$h'$, refining the partitions if necessary:
\[
 \ell(h')\le\ell(\bar h_m)
       -\frac12\int_0^1(h'-\bar h_m)\rho_{\bar h_m}.
\]
The fields $\bar h_m$ are uniformly bounded and converge to $h$ in
$L^1$.  Using \eqref{cav:rho-convergence} and continuity of $\ell$, we
may pass to the limit to get
\[
 \ell(h')+\frac12\int_0^1h'p
 \le\ell(h)+\frac12\int_0^1hp.
\]
Taking the supremum over all allowed field steps gives the upper bound
in \eqref{cav:entropy-identity}.  Conversely, approximate the bounded
nondecreasing $h$ in $L^1$ by nonnegative nondecreasing step functions.
Continuity of $\ell$ and boundedness of $p$ show that their dual values
converge to the right side of \eqref{cav:entropy-identity}.  The
definition of $S_{\mathrm I}$ gives the reverse inequality.
\end{proof}

\subsection{Uniform increments and telescoping}

\begin{proposition}\label{cav:lower-bound}
For $f\in C_c^\infty(\mathbb R)$,
\[
 \liminf_{N\to\infty}\mathbb E p_N
 \ge\inf_q\{S_{\mathrm I}(q)+\alpha V_f(q)\}.
\]
\end{proposition}

\begin{proof}
Write $\mathcal P=\inf_q\{S_{\mathrm I}(q)+\alpha V_f(q)\}$, a finite
number by the elementary bounds for the variational functional.  The
limit \eqref{cav:increment-limit} and
Proposition~\ref{cav:dual-attainment} give, on every subsequence
constructed above,
\[
 \lim_N\Delta_{N,L}(v(N))
 =L\{S_{\mathrm I}(p)+\alpha V_f(p)\}
                +(d-\alpha L)V_f(p)
 \ge L\mathcal P-\|f\|_\infty.
\]
We used $|V_f(p)|\le\|f\|_\infty$ and $|d-\alpha L|\le1$.
Crucially, the starting good parameter sequence was arbitrary.
Consequently
\begin{equation}\label{cav:uniform-liminf}
 \liminf_{N\to\infty}\inf_{v\in\mathcal G_N}\Delta_{N,L}(v)
 \ge L\mathcal P-\|f\|_\infty.
\end{equation}
For otherwise a fixed positive violation could be selected at
successive dimensions with deterministic $v(N)\in\mathcal G_N$;
compactness would supply the array and constant-$d$ subsequence above,
contradicting its evaluated limit.

The uniform all-parameter lower bound of
Lemma~\ref{cav:finite-comparison} controls the complement of
$\mathcal G_N$.  Specifically, for every $\varepsilon>0$ and all large
$N$, the parameter integral is at least
\[
 \pi(\mathcal G_N)
       (L\mathcal P-\|f\|_\infty-\varepsilon)
       -C_L\pi(\mathcal G_N^c).
\]
The use of the same product probability in both dimensions gives the
exact identity
\[
 \int\Delta_{N,L}(v)\,\pi(dv)=\mathcal F_{N+L}-\mathcal F_N.
\]
No good-parameter condition in dimension $N+L$ is imposed or needed.
We conclude that
\[
 \liminf_N(\mathcal F_{N+L}-\mathcal F_N)
       \ge L\mathcal P-\|f\|_\infty.
\]
For fixed $L$ and $\varepsilon>0$, telescope the eventual bound on
these increments separately in each of the $L$ residue classes modulo
$L$.  The finitely many initial terms in each class vanish upon division
by $N$, so
\[
 \liminf_N\frac{\mathcal F_N}{N}
 \ge\mathcal P-\frac{\|f\|_\infty+\varepsilon}{L}.
\]
First let $\varepsilon\downarrow0$ and then let $L\to\infty$.
Finally \eqref{bulk:perturbation-cost} transfers the result to the
unperturbed pressure.  Together with the upper bound this proves the
variational formula for smooth compactly supported $f$.
\end{proof}
